% =============================================================================
% GitHub Copilot — Token Optimization Under the New Billing
% Beamer deck (Metropolis theme), 16:9
%
% Compile:
%   latexmk -pdf -xelatex copilot-token-optimization.tex
% or
%   xelatex copilot-token-optimization.tex   (run twice for ToC + page numbers)
%
% pdflatex also works but Metropolis renders best with XeLaTeX or LuaLaTeX
% because of the Fira Sans font.
% =============================================================================
\documentclass[11pt,aspectratio=169]{beamer}

% --- Theme -------------------------------------------------------------------
\usetheme{metropolis}
\metroset{
  progressbar=frametitle,
  numbering=fraction,
  block=fill,
  sectionpage=progressbar,
  subsectionpage=progressbar
}

% --- Encoding & fonts --------------------------------------------------------
\usepackage[english]{babel}
\usepackage{microtype}
\usepackage{fontspec}            % Metropolis prefers XeLaTeX/LuaLaTeX

% --- Colour palette (overrides) ----------------------------------------------
\definecolor{accentblue}{HTML}{2563EB}
\definecolor{darkgray}{HTML}{1F2937}
\definecolor{lightgray}{HTML}{F3F4F6}
\definecolor{successgreen}{HTML}{059669}
\definecolor{warningred}{HTML}{DC2626}
\setbeamercolor{progress bar}{fg=accentblue}
\setbeamercolor{frametitle}{bg=darkgray}
\setbeamercolor{title separator}{fg=accentblue}

% --- Tables, graphics, code --------------------------------------------------
\usepackage{booktabs}
\usepackage{array}
\usepackage{graphicx}
\usepackage{tikz}
\usepackage{listings}
\usepackage{xcolor}
\lstdefinestyle{prompt}{
  basicstyle=\ttfamily\footnotesize,
  backgroundcolor=\color{lightgray},
  frame=none,
  xleftmargin=10pt,
  framexleftmargin=10pt,
  breaklines=true,
  showstringspaces=false,
  literate={→}{{$\rightarrow$}}1
}
\lstset{style=prompt}

% --- Refs --------------------------------------------------------------------
% hyperref is auto-loaded by beamer; just configure it here to avoid option clash
\hypersetup{hidelinks}

% --- Custom commands ---------------------------------------------------------
\newcommand{\good}{\textcolor{successgreen}{\checkmark}}
\newcommand{\bad}{\textcolor{warningred}{\textbf{\texttimes}}}
\newcommand{\hl}[1]{\textbf{\textcolor{accentblue}{#1}}}

% --- Title -------------------------------------------------------------------
\title{GitHub Copilot}
\subtitle{Token Optimization Under the New Billing}
\author{Internal session}
\institute{Office training pack}
\date{June 2026}

% =============================================================================
\begin{document}

\maketitle

% =============================================================================
\begin{frame}{Agenda}
  \begin{itemize}
    \item What changed on 1 June 2026
    \item The five VS Code Copilot surfaces and what each costs
    \item Four cost levers, in priority order
    \item The two-lane workflow
    \item Quick wins to adopt this week
  \end{itemize}
  \vfill
  \small Time: 20 minutes plus Q\&A.
\end{frame}

% =============================================================================
\section{What changed}

\begin{frame}{What changed on 1 June 2026}
  \begin{itemize}
    \item \textbf{Old system:} one chat turn $=$ one ``premium request'', multiplied by a per-model factor. Plans gave us \hl{N requests per month}.
    \item \textbf{New system:} plans give us a \hl{dollar budget in AI Credits}. Every paid surface bills against it \hl{by tokens $\times$ per-model rate}.
  \end{itemize}
  \vfill
  \begin{block}{Three things to remember}
    \begin{itemize}
      \item \hl{1 AI Credit $=$ \$0.01}
      \item Code completion and Next Edit Suggestions remain \hl{free}
      \item Chat, Edit mode, Agent mode, Code Review are \hl{paid}
    \end{itemize}
  \end{block}
\end{frame}

\begin{frame}{Why this matters}
  A realistic Agent-mode session on Opus can spend \hl{\$0.50 to \$5+} in credits.

  \medskip
  Public reports document developers burning \hl{\$30--40 in a single morning} on long agentic sessions.

  \medskip
  Pro is \hl{\$10/month}. The math is unforgiving if we are not deliberate.

  \vfill
  \begin{alertblock}{The good news}
    The highest-frequency Copilot interactions --- inline completion and Next Edit Suggestions --- remain \hl{unmetered}. Most of the productivity we get from Copilot is in the free lane.
  \end{alertblock}
\end{frame}

\begin{frame}{Plan allowances}
  \centering
  \begin{tabular}{lrr}
    \toprule
    \textbf{Plan} & \textbf{Monthly cost} & \textbf{Monthly AI Credits} \\
    \midrule
    Copilot Pro            & \$10        & \$10 \\
    Copilot Pro+           & \$39        & \$39 \\
    Copilot Business       & \$19 / seat & \$19 / seat \\
    Copilot Enterprise     & \$39 / seat & \$39 / seat \\
    \bottomrule
  \end{tabular}

  \vfill
  \small No rollover. Overage billable per token (admin policy applies).
\end{frame}

% =============================================================================
\section{The surfaces}

\begin{frame}{The five surfaces in VS Code}
  \centering\small
  \begin{tabular}{lcl}
    \toprule
    \textbf{Surface} & \textbf{Bills?} & \textbf{Typical cost shape} \\
    \midrule
    Code completion (ghost text)        & \good\ no  & Free \\
    Next Edit Suggestions               & \good\ no  & Free \\
    Chat --- Ask                        & \bad\ yes  & Small per turn \\
    Chat --- Edit                       & \bad\ yes  & Medium --- whole file in \& out \\
    Chat --- Agent                      & \bad\ yes  & \hl{Largest} --- tool output stacks \\
    Code Review (PR / inline)           & \bad\ yes  & Per review, $\propto$ diff size \\
    \bottomrule
  \end{tabular}
\end{frame}

\begin{frame}[fragile]{Where the tokens come from}
  On every paid turn:
\begin{lstlisting}
cost = (prompt + attached context + chat history
        + tool output history)      x  input_rate
     + (model reply + tool calls)   x  output_rate
\end{lstlisting}
  \medskip
  Four controllable levers, in priority order:
  \begin{enumerate}
    \item \hl{Which model} --- 10$\times$ spread between cheapest and most expensive
    \item \hl{Tool output history} --- agent mode only, but dominates totals there
    \item \hl{Attached context size} --- \texttt{\#selection} vs \texttt{@workspace} vs dragged folder
    \item \hl{Chat history length} --- new chat per topic
  \end{enumerate}
\end{frame}

% =============================================================================
\section{The levers}

\begin{frame}{Lever 1 --- Model selection}
  The \hl{single highest-leverage decision} you make per turn.

  \medskip
  \centering\small
  \begin{tabular}{lll}
    \toprule
    \textbf{Tier} & \textbf{Examples} & \textbf{When to use} \\
    \midrule
    Cheap    & Haiku 4.5, GPT-5 mini, Gemini 3 Flash & Exploration, summarising, simple edits \\
    Standard & Sonnet 4.6, GPT-5.4, Gemini 3.1 Pro   & Multi-file edits, careful implementation \\
    Premium  & Opus 4.7, GPT-5.5                     & Hard reasoning, debugging, architecture \\
    \bottomrule
  \end{tabular}

  \vfill
  \raggedright
  \textbf{Action:} in VS Code, set your default chat model to a \hl{cheap-lane} model. Escalate deliberately, not reflexively.
\end{frame}

\begin{frame}[fragile]{The two-lane workflow}
\begin{lstlisting}
EXPLORE  ->  Cheap lane (Haiku, GPT-5 mini)
             ask . plan . summarise . iterate freely
                          |
                          v
EXECUTE  ->  Premium lane (Sonnet+ / Opus)
             one well-formed implementation turn
                          |
                          v
VERIFY   ->  Cheap lane again
             review the diff . write tests
\end{lstlisting}
  \vfill
  Framing: \hl{expensive models for executing locked decisions; cheap models for exploring undecided ones}.
\end{frame}

\begin{frame}[fragile]{Lever 2 --- Context size}
  Every attached token is paid input \hl{on every turn it stays in the conversation}.

  \medskip
  Cost ranking, cheapest to most expensive:
\begin{lstlisting}
#selection  <  #editor  <  #file:<name>  <  @workspace  <  dragged folder
\end{lstlisting}

  \medskip
  Rules of thumb:
  \begin{itemize}
    \item 20 lines matter? \texttt{\#selection} them.
    \item One file matters? \texttt{\#file:} it --- not \texttt{@workspace}.
    \item \texttt{@workspace} is for ``I don't know which files matter yet.''
    \item \hl{Never drag a folder} unless that is truly the intent.
  \end{itemize}
\end{frame}

\begin{frame}[fragile]{Lever 3 --- Conversation hygiene}
  \textbf{Anti-pattern --- prompt-chipping:}
\begin{lstlisting}
"Refactor this to use a dict."
"Actually use defaultdict."
"Make the keys strings."
"Add type hints."
\end{lstlisting}
  Four turns, four bills, each carrying the whole history.

  \medskip
  \textbf{Pattern --- one complete brief:}
\begin{lstlisting}
"Refactor this to use a
 collections.defaultdict[str, list[Foo]],
 add type hints, keep the existing public signature."
\end{lstlisting}
  One turn. One bill. Often better output too.
\end{frame}

\begin{frame}{Lever 4 --- Agent mode discipline}
  Agent mode is where the budget vanishes if you are casual.

  \medskip
  \textbf{Front-load the brief:}
  \begin{itemize}
    \item \textbf{Task} --- one sentence
    \item \textbf{In scope / Out of scope} --- which files; what NOT to touch
    \item \textbf{Constraints} --- patterns to preserve, libs required/forbidden
    \item \textbf{Definition of Done} --- how the agent knows to stop
    \item \textbf{Reporting} --- summarise files touched and assumptions
  \end{itemize}

  \vfill
  A well-shaped brief routinely \hl{cuts agent sessions from 20 turns to 4--6}.
\end{frame}

\begin{frame}{Watch the token counter}
  The chat panel shows the running token usage in the \hl{bottom-right corner}.

  \medskip
  \centering
  \begin{tabular}{ll}
    \toprule
    \textbf{Tokens in context} & \textbf{What to do} \\
    \midrule
    Under 30K & Comfortable, continue \\
    30--80K   & Expensive on premium --- check whether you are going in circles \\
    80K+      & Stop. Summarise to a fresh chat. \\
    \bottomrule
  \end{tabular}
\end{frame}

% =============================================================================
\section{Reusable context \& habits}

\begin{frame}[fragile]{Reusable context lives in files}
  Stop re-pasting ``remember, we follow X convention''. Put it in a file Copilot auto-loads.

  \medskip
  At repo root --- \texttt{.github/copilot-instructions.md}:
\begin{lstlisting}
- We use Python 3.12, type hints required.
- No bare `except`. Use specific exceptions.
- SQLAlchemy 2.x -- no Query API; use select().
- Tests in `tests/`, one file per source module.
\end{lstlisting}

  \medskip
  For repeated task shapes --- \texttt{.github/prompts/<name>.prompt.md}:
  reusable templates for ``write tests in our style'', ``refactor to our pattern''.
\end{frame}

\begin{frame}{Quick wins --- adopt this week}
  \begin{itemize}
    \item \good\ \textbf{Default model} $\rightarrow$ Haiku 4.5 or GPT-5 mini in VS Code
    \item \good\ \textbf{Bind a shortcut} $\rightarrow$ ``New Chat'' for one-keystroke restart
    \item \good\ \textbf{\texttt{.github/copilot-instructions.md}} in every active repo
    \item \good\ \textbf{Stop drag-folder} entirely; always \texttt{\#file:} or \texttt{@workspace}
    \item \good\ \textbf{One brief, not four chips} --- write the full prompt
    \item \good\ \textbf{Glance at the token counter} before big turns
  \end{itemize}
\end{frame}

\begin{frame}{What we are trading}
  These habits feel like upfront work.

  \medskip
  The trade is \hl{upfront thinking} vs.\ \hl{downstream burn}:
  \medskip
  \centering
  \begin{tabular}{lr}
    \toprule
    \textbf{Habit} & \textbf{Cost} \\
    \midrule
    Writing a complete prompt        & +1 min thinking \\
    Switching models mid-task        & +5 sec clicks \\
    Starting a new chat per topic    & +2 sec \\
    \midrule
    \textbf{Effect on credit burn}   & \hl{$-50\%$ to $-80\%$} typical \\
    \textbf{Effect on output quality}& Usually meaningfully better \\
    \textbf{Effect on rework}        & Lower \\
    \bottomrule
  \end{tabular}
\end{frame}

% =============================================================================
\section{Summary}

\begin{frame}{Summary --- the five rules}
  \begin{enumerate}
    \item \hl{Free is free} --- completion and NES don't bill. Default to them whenever the task is shaped like ``complete this''.
    \item \hl{Cheap model first} --- explore on Haiku or mini; escalate only to execute a locked plan.
    \item \hl{Attach less} --- \texttt{\#selection} before \texttt{\#editor} before \texttt{\#file:} before \texttt{@workspace}.
    \item \hl{One task $=$ one chat} --- new chat per topic change.
    \item \hl{Agent mode $=$ high-stakes mode} --- front-load the brief, set the Definition of Done.
  \end{enumerate}
\end{frame}

\begin{frame}{Resources we maintain}
  Everything in this talk is documented in the team folder:
  \begin{itemize}
    \item \texttt{README.md} --- overview and how to use the pack
    \item \texttt{01-new-billing-model.md} --- the billing system, allowances, model rate card
    \item \texttt{02-vscode-surfaces.md} --- per-surface deep dive
    \item \texttt{03-optimization-playbook.md} --- the full playbook
    \item \texttt{04-cheatsheet.md} --- one-page printable reference
    \item \texttt{appendix-sources.md} --- every claim, where it came from
  \end{itemize}
\end{frame}

% =============================================================================
\begin{frame}[plain]
  \begin{center}
    {\Huge Thank you}\\[1.2em]
    {\Large Questions?}\\[2.4em]
    \small \textit{Verified June 2026 against docs.github.com/en/copilot/reference/copilot-billing.}
  \end{center}
\end{frame}

\end{document}
